%!TEX root=../mythesis.tex
% Chapter 3 -- EarnHFT

\chapter{EarnHFT: Hierarchical Reinforcement Learning for High-Frequency Trading} % Main chapter title
\chaptermark{EarnHFT for High-Frequency Trading}
\label{ch:earnhft}

%----------------------------------------------------------------------------------------
\section{Introduction}
\label{sec:earnhft_intro}

High-frequency trading (HFT) uses computer algorithms to place and cancel orders at very short time scales, and accounts for a majority of volume in modern financial markets~\cite{ALMEIDA2023100785}. It is especially prevalent in the cryptocurrency market, whose round-the-clock trading removes overnight risk and whose dramatic price fluctuations create frequent short-lived opportunities. A profitable HFT strategy can outperform a low-frequency one, which is why it is pursued aggressively by quantitative traders.

Reinforcement learning has produced strong results on low-frequency tasks such as daily stock or futures trading~\cite{sun2022deepscalper,briola2021deep}, but these methods rarely survive the move to HFT, for two reasons identified in \sref{sec:bg_gaps} as Gap~I.

\begin{itemize}
\item \textbf{The horizon is extreme, so RL is data-inefficient.} A second-level agent must be evaluated over dozens of days, giving a horizon of roughly one million steps---for instance, $60\times 60\times 24\times 12 = 1{,}036{,}800$ seconds over twelve days---against the few-thousand-step horizon of Atari benchmarks~\cite{mnih2013playing}. Such horizons demand far more data to converge~\cite{zhang2023towards}, straining computation.
\item \textbf{The market drifts, so a single agent cannot maintain performance.} Standard RL assumes the training and testing environments coincide. In cryptocurrency, the market trend changes substantially between them; an agent trained under one trend incurs large losses when the trend reverses.
\end{itemize}

This chapter presents EarnHFT, an \textbf{E}fficient hier\textbf{A}rchical \textbf{R}einforcement lear\textbf{N}ing method for \textbf{H}igh \textbf{F}requency \textbf{T}rading, whose three stages are summarised in \fref{fig:earnhft_pipeline}. The design follows the decompose--specialise--select pattern of \sref{sec:intro_pattern}.

\begin{enumerate}
\item \textbf{Stage~I (Q-teacher).} We exploit a structural property of execution---that at second scale a single trader's orders do not move prices, so future prices are independent of the policy---to compute, by dynamic programming, the optimal action value for every state. Used as a regulariser and a demonstrator, this ``Q-teacher'' accelerates and stabilises second-level training.
\item \textbf{Stage~II (specialist pool).} We train hundreds of agents under different market-trend preferences, segment and label the market without hyperparameter tuning, and select a small pool of the most profitable specialists.
\item \textbf{Stage~III (router).} We train a minute-level router that dynamically picks a second-level specialist from the pool to maintain performance as the market changes.
\end{enumerate}

Across four cryptocurrency markets spanning bull, bear, and sideways conditions, EarnHFT outperforms six state-of-the-art baselines on six financial metrics, exceeding the runner-up by at least 30\% in profitability.

\begin{figure}[!htb]
\centering
\includegraphics[width=\textwidth]{Chapter4/earnhft_pipeline}
\caption[The EarnHFT framework.]{Overview of the three stages of EarnHFT. Stage~I computes a Q-teacher to enhance the performance and training efficiency of second-level agents. Stage~II trains diverse agents under various market trends and selects a small fraction by profitability to form an agent pool. Stage~III trains a minute-level router that dynamically picks a second-level agent from the pool.}
\label{fig:earnhft_pipeline}
\end{figure}

%----------------------------------------------------------------------------------------
\section{Problem Formulation}
\label{sec:earnhft_formulation}

The financial concepts---limit order book, OHLC, technical indicators, market order, position, and net value---are defined in \sref{sec:bg_instruments}. We restrict attention to a single asset with long-only positions, and the objective is to maximise net value $V_t=V_{ct}+H_t\, p_t^{b_1}$, where $V_{ct}$ is cash and $H_t$ the position valued at the best bid. The execution price of a market order walks the book, $E_t(M)=\sum_i (p_t^i\min(q_t^i,R_{i-1}))(1+\sigma)$, where $R_{i-1}$ is the residual quantity after level $i$ and $\sigma$ the commission rate.

\subsection{A Hierarchical Markov Decision Process}
\label{sec:earnhft_hmdp}

The two challenges of \sref{sec:earnhft_intro} have a common structural cause: micro-level information drifts, so no single agent profits under all trends, while macro-level information, an aggregation of the micro level, carries the slow trend signal. We therefore decompose HFT into a hierarchical MDP $(\mathrm{MDP}_h,\mathrm{MDP}_l)$, in which a low-level MDP at second scale models micro-level dynamics and execution, and a high-level MDP at minute scale models macro-level trend and strategy switching:
\begin{align*}
\mathrm{MDP}_h &= (S_h,A_h,P_h,R_h,\gamma_h,T_h), &
\mathrm{MDP}_l &= (S_l,A_l,P_l,R_l,\gamma_l,T_l).
\end{align*}

\paragraph{Low level (second scale).} The state $S_{lt}$ comprises a latent representation $y_{lt}$ of 54 micro-level technical indicators computed from a length-60 rolling window of second-level OHLC and LOB snapshots, together with the current position $H_t$. The action $a_{lt}$ is a target position chosen from a finite pool $\{0,\frac{H}{|A|-1},\dots,H\}$; if it differs from the current position, a market order is executed immediately. The reward $r_{lt}=H_{t+1}p_{t+1}^{b_1}-(H_t p_t^{b_1}+E_t(H_{t+1}-H_t))$ is the per-second change in net value.

\paragraph{High level (minute scale).} The state $S_{ht}$ comprises a latent representation $y_{ht}$ of 19 macro-level indicators from a length-60 rolling window of minute-level OHLC, together with $H_t$. The action $a_{ht}$ selects an agent from the pre-trained pool $A_h$. The reward $r_{hT}=\sum_{t=T}^{T+\tau}r_{lt}$ is the net-value change accrued by the selected low-level agent over one minute. In summary, every minute the high-level router picks a low-level agent, which then adjusts its position every second; the goal is to find a set of specialists and a router that jointly maximise total profit.

%----------------------------------------------------------------------------------------
\section{Method}
\label{sec:earnhft_method}

\subsection{Stage I: Efficient Reinforcement Learning with a Q-Teacher}
\label{sec:earnhft_stage1}

A long trajectory normally inflates the cost of RL. In the low-level MDP, however, prices are not influenced by the policy, so future prices are known in hindsight on historical data. This lets us construct the optimal action value $Q^\star$ by dynamic programming~\cite{sutton2018reinforcement} and use it to teach the agent, both as a value supervisor and as a demonstrator.

\paragraph{Optimal value supervisor.} Because the position set is finite, the optimal action value can be computed backwards: from the last step, the value of each (position, action) pair is the immediate reward plus the best continuation value, as in Algorithm~\ref{alg:earnhft_qtable}. We add to the DDQN loss a supervision term that is the Kullback--Leibler divergence between the agent's action values and the optimal action values at the same state:
\begin{equation}
\label{eq:earnhft_loss}
L(\theta_i) = L_{td} + \alpha\,\mathrm{KL}\!\left(Q_t(\chi,p,\cdot;\theta_i)\,\|\,Q^\star(\chi,p,\cdot)\right),
\end{equation}
where $L_{td}=\left(r+\gamma\max Q_t(\chi',a,\cdot;\theta_i')-Q_t(\chi,p,a;\theta_i)\right)^2$ is the DDQN temporal-difference error, $\chi$ the latent representation, $p$ the position, and $\alpha$ a coefficient decaying over time. The supervision term lets the agent learn the relative advantage of all actions at a state without exploring them, accelerating training.

\begin{algorithm}[!htb]
\caption{Construction of the Optimal Action Value}
\label{alg:earnhft_qtable}
\begin{algorithmic}[1]
\State \textbf{Input:} time series $\mathcal{D}$ of length $N$, commission rate $\delta$, action space $A$
\State \textbf{Output:} table $Q^\star$ of optimal action values indexed by time, position, action
\State Initialise $Q^\star$ with shape $(N,|A|,|A|)$ and all elements $0$.
\For{$t \gets N-1$ \textbf{to} $1$}
  \For{$p \gets 1$ \textbf{to} $|A|$}
    \For{$a \gets 1$ \textbf{to} $|A|$}
      \State $Q^\star[t,p,a] \gets \max_{a'} Q^\star[t+1,a,a'] + a\,p_{t+1}^{b_1} - (p\,p_{t}^{b_1} + E_t(p-a))$
    \EndFor
  \EndFor
\EndFor
\State \textbf{return} $Q^\star$
\end{algorithmic}
\end{algorithm}

\paragraph{Optimal actor.} Supervision alone is insufficient, because the optimal action value is conditioned on the current position: once the agent deviates from the optimal policy, the supervision target shifts, compounding the deviation. We therefore additionally train on transitions generated by the optimal policy itself---the action with the highest optimal value---which supplies demonstration experience and prevents the agent from settling in a local trap. Algorithm~\ref{alg:earnhft_efficient} collects experience from both the $\epsilon$-greedy policy and the optimal actor, then updates the network using both TD errors and the KL supervision.

\begin{algorithm}[!htb]
\caption{Efficient Reinforcement Learning with a Q-Teacher}
\label{alg:earnhft_efficient}
\begin{algorithmic}[1]
\State \textbf{Input:} time series $\mathcal{D}$ of length $N$, commission rate $\delta$, action space $A$
\State \textbf{Output:} network parameters $\theta$
\State Initialise replay buffer $R$, network $Q_\theta$, target $Q_{\theta'}$; build $Q^\star$ via Algorithm~\ref{alg:earnhft_qtable}; initialise environment.
\For{$t = 1$ \textbf{to} $N-1$} \Comment{exploration rollout}
  \State Choose $a_\epsilon$ by the $\epsilon$-greedy policy; store $(s,a_\epsilon,r,s',Q^\star)$ in $R$.
\EndFor
\For{$t = 1$ \textbf{to} $N-1$} \Comment{optimal-actor rollout}
  \State Choose $a_o = \argmax_{a} Q^\star[t,p,a]$; store $(s,a_o,r,s',Q^\star)$ in $R$.
\EndFor
\State Sample transitions; compute $L$ by \eqref{eq:earnhft_loss}; descend on $\theta$ and set $\theta'=\tau\theta+(1-\tau)\theta'$.
\State \textbf{return} $Q_\theta$
\end{algorithmic}
\end{algorithm}

The Q-teacher does not compromise correctness. Treating the supervised update as an operator $H$ on action values, one can show $H$ is a contraction in the sup-norm whose fixed point is $Q^\star$, so the learned value still converges to the optimum. Following the convergence argument for Q-learning~\cite{jaakkola1993convergence}, it suffices to verify these two properties of $H$. The operator combines the usual Bellman backup with the supervision term and is
\begin{equation}
\label{eq:earnhft_operator}
Hq(x,a) = \lambda\sum_{y\in X}T(y\mid x,a)\bigl(r(x,a,y)+\gamma\max_b q(y,b)\bigr) + (1-\lambda)\Bigl(\tfrac{1}{|A|}\textstyle\sum_{a\in A}q(x,a)+Re(x,a)\Bigr),
\end{equation}
where the supervision residual is $Re(x,a) = q^\star(x,a) - \tfrac{1}{|A|}\sum_{a\in A}q^\star(x,a)$ and $\lambda\in(0,1]$ weights the Bellman backup against the supervision term.

\begin{proposition}[$Q^\star$ is a fixed point of $H$]
\label{prop:earnhft_fixed}
The optimal action value $q^\star$ satisfies $Hq^\star = q^\star$.
\end{proposition}
\begin{proof}
Substituting $q^\star$ into \eqref{eq:earnhft_operator} and using the Bellman optimality equation $q^\star(x,a)=\sum_y T(y\mid x,a)(r(x,a,y)+\gamma\max_b q^\star(y,b))$ for the first term and the definition of $Re$ for the second,
\begin{align*}
Hq^\star(x,a) &= \lambda\, q^\star(x,a) + (1-\lambda)\Bigl(\tfrac{1}{|A|}\textstyle\sum_a q^\star(x,a) + q^\star(x,a) - \tfrac{1}{|A|}\textstyle\sum_a q^\star(x,a)\Bigr) \\
&= \lambda\, q^\star(x,a) + (1-\lambda)\, q^\star(x,a) = q^\star(x,a). \qedhere
\end{align*}
\end{proof}

\begin{proposition}[Contraction in the sup-norm]
\label{prop:earnhft_contraction}
For any action values $q_1,q_2$, $\|Hq_1-Hq_2\|_\infty \le k\,\|q_1-q_2\|_\infty$ with $k=1-(1-\lambda)\gamma<1$.
\end{proposition}
\begin{proof}
The supervision residual $Re(x,a)$ depends only on $q^\star$ and not on the argument of $H$, so it cancels in the difference $Hq_1-Hq_2$. Bounding the $\max$ operator by $|\max_b q_1(y,b)-\max_b q_2(y,b)|\le\max_b|q_1(y,b)-q_2(y,b)|$,
\begin{align*}
\|Hq_1-Hq_2\|_\infty
&\le \lambda\gamma\,\|q_1-q_2\|_\infty + (1-\lambda)\,\|q_1-q_2\|_\infty \\
&= \bigl(1-(1-\lambda)\gamma\bigr)\|q_1-q_2\|_\infty.
\end{align*}
Since $\gamma\in(0,1]$ and $\lambda\in(0,1]$, the factor $k=1-(1-\lambda)\gamma<1$.
\end{proof}

\begin{theorem}[Convergence of the supervised update]
\label{thm:earnhft_convergence}
Under a finite MDP, the action value learned with the optimal-value-supervisor update of \eqref{eq:earnhft_loss} converges to the optimal action value $Q^\star$.
\end{theorem}
\begin{proof}
Propositions~\ref{prop:earnhft_fixed} and~\ref{prop:earnhft_contraction} show that $H$ is a sup-norm contraction with unique fixed point $Q^\star$. By the stochastic-approximation argument of~\cite{jaakkola1993convergence}, the iterates therefore converge almost surely to $Q^\star$.
\end{proof}

\subsection{Stage II: Construction of the Agent Pool}
\label{sec:earnhft_stage2}

Because micro-level information changes rapidly, no single agent maintains performance over a long period; preliminary experiments confirm that training across conflicting trends is incompatible. We therefore decompose the market into trends and develop a specialist for each.

\paragraph{Generating diverse agents.} Prior approaches obtain diversity as a by-product of random seeds or hyperparameter search~\cite{sun2023mastering}, which is unstructured. We instead train agents with explicit \emph{preferences} over market trends. The training series (over three million steps) is cut into chunks of length $L$, each a continuous trend; a preference parameter $\beta$ assigns each chunk with buy-and-hold return rate $r$ a sampling priority
\begin{equation}
\label{eq:earnhft_pes}
f(x) =
\begin{cases}
\dfrac{e^{\beta r}}{\mathrm{pdf}(r)} & \text{if } Q_{\frac{\theta}{2}}(R) \le r \le Q_{1-\frac{\theta}{2}}(R),\\[2mm]
e^{\beta r} & \text{otherwise},
\end{cases}
\end{equation}
where $Q_{\frac{\theta}{2}}(R)$ is the $\frac{\theta}{2}$-quantile of return rates and $\mathrm{pdf}$ is a kernel density estimate,
\begin{equation}
\label{eq:earnhft_kernel}
\mathrm{pdf}(x)=\frac{1}{nh}\sum_{r\in R} K\!\left(\frac{x-r}{h}\right),
\end{equation}
with bandwidth $h$ searched around Silverman's rule~\cite{silverman1984spline} and $K$ a standard normal kernel. Dividing by the density removes the influence of the dataset's own distribution, while $e^{\beta r}$ tilts sampling towards a preferred trend; clipping the extremes prevents a few violent chunks from poisoning the agent everywhere. Different $\beta$ values yield agents specialised to different trends.

\paragraph{Market segmentation and labelling.} To evaluate specialists per trend we need to label the market. Algorithm~\ref{alg:earnhft_segment} denoises the series, splits it at extrema, merges adjacent segments whose dynamic-time-warping~\cite{muda2010voice} and slope differences are small, and labels segments by slope quantiles. Unlike prior methods, it requires no per-dataset hyperparameter tuning.

\begin{algorithm}[!htb]
\caption{Market Segmentation and Labelling}
\label{alg:earnhft_segment}
\begin{algorithmic}[1]
\State \textbf{Input:} time series $\mathcal{D}$ of length $N$; risk threshold $\theta$; number of labels $M$
\State \textbf{Output:} a trend label for every point of $\mathcal{D}$
\State $\mathcal{D}' \gets$ denoise high-frequency noise of $\mathcal{D}$.
\State Divide $\mathcal{D}'$ at its extrema into segments $S$.
\State Merge adjacent segments in $S$ whose DTW and slope differences are small, until $S$ is stable.
\State Compute thresholds $H = Q_{1-\frac{\theta}{2}}(R)$, $L = Q_{\frac{\theta}{2}}(R)$ on segment slopes $R$.
\State Compute per-label slope bounds from the quantiles and thresholds; label each segment.
\State \textbf{return} the label of each segment.
\end{algorithmic}
\end{algorithm}

\paragraph{Agent selection.} Including all trained agents would inflate the router's action space, so we select a small pool. Each point of the validation set is labelled by Algorithm~\ref{alg:earnhft_segment}; agents are evaluated across trends and initial positions, and the most profitable agent under each (trend, initial-position) pair is retained, forming a two-dimensional pool indexed by $(m,n)$ for $m$ trends and $n$ initial positions.

\subsection{Stage III: Dynamic Routing}
\label{sec:earnhft_stage3}

We train the router for the high-level MDP with DDQN~\cite{van2016deep}. Operating at minute scale shortens the trajectory by 98.33\% relative to the second scale, and we further restrict the router, at each step, to agents whose initial position matches the current position, reducing the choice to $m$ agents. We deliberately do \emph{not} compute a Q-teacher for the router: the horizon is already short, and the high-level action is a strategy rather than a target position, which would make the optimal value expensive to compute. With the computation for self-exploration reduced and that for the teacher increased, plain DDQN is the more efficient choice at the high level.

%----------------------------------------------------------------------------------------
\section{Experimental Setup}
\label{sec:earnhft_setup}

\paragraph{Datasets.} We evaluate on four cryptocurrencies covering mainstream and niche assets and bull, bear, and sideways conditions, summarised in \tref{tab:earnhft_dataset}. On each, the last nine days are used for testing, the preceding nine for validation, and the remainder for training. Low-level agents are trained on the training set; the validation set is segmented and labelled for model selection; the router is then trained on the training set and selected on the full validation set before testing.

\begin{table}[!htb]
\setlength\tabcolsep{6pt}
\centering
\caption[EarnHFT dataset statistics.]{The four cryptocurrency datasets, with market dynamics, number of seconds, and chronological span. Dates are in YY/MM/DD format.}
\label{tab:earnhft_dataset}
\begin{tabular}{lcccc}
\toprule
\textbf{Dataset} & \textbf{Dynamics} & \textbf{Seconds} & \textbf{From} & \textbf{To} \\
\midrule
BTC/TUSD & Sideways & 4{,}057{,}140 & 23/03/30 & 23/05/15 \\
BTC/USDT & Sideways & 3{,}884{,}400 & 22/09/01 & 22/10/15 \\
ETH/USDT & Bear     & 3{,}970{,}800 & 22/05/01 & 22/06/15 \\
GALA/USDT & Bull    & 3{,}970{,}740 & 22/07/01 & 22/08/15 \\
\bottomrule
\end{tabular}
\end{table}

\paragraph{Metrics.} We report the six financial metrics defined in \sref{sec:bg_metrics}: TR (profit); ASR, ACR, ASoR (risk-adjusted profit); and AVOL, MDD (risk).

\paragraph{Baselines.} We compare against six methods: two policy-based RL algorithms, PPO~\cite{schulman2017proximal} and DRA~\cite{briola2021deep}; two value-based RL algorithms, DQN~\cite{mnih2015human} and CDQNRP~\cite{zhu2022quantitative}; and two rule-based methods, MACD~\cite{krug2022enforcing} and IV~\cite{chordia2002order}. Where official implementations exist we reuse their hyperparameters; otherwise we reimplement faithfully from the original papers.

\paragraph{Training.} All experiments run on a single RTX~4090 GPU. The commission rate is $0$ for BTC/TUSD and $0.02\%$ for the others, following Binance. We train with $\beta\in\{-90,-10,30,100\}$ for 50 epochs each, producing 200 agents, optimised with Adam. The full set of experiments completes in about ten hours.

%----------------------------------------------------------------------------------------
\section{Results and Analysis}
\label{sec:earnhft_results}

\subsection{Comparison with Baselines}
\label{sec:earnhft_comparison}

\tref{tab:earnhft_main} reports the main comparison. EarnHFT achieves the highest profit (TR) on all four datasets and the highest risk-adjusted profit on three. Value-based methods (DQN, CDQNRP) perform acceptably only when the validation and test sets are similar under a stable trend. Policy-based methods (PPO, DRA) tend to collapse onto a do-nothing policy---returning a target equal to the current position to avoid commission---even at a learning rate of $10^{-7}$, and fail badly in the bear market. Rule-based methods are highly sensitive to their take-profit and stop-loss thresholds and earn only moderate profit in volatile markets. EarnHFT, by contrast, is a deliberately aggressive trader because the optimal-value supervisor and optimal actor convey profit-oriented experience; it consequently trails on some risk metrics while dominating on profit. \fref{fig:earnhft_trading} shows a representative episode: EarnHFT opens and closes a position within thirty seconds, profiting from a movement that appears as a mere pullback at minute scale.

\begin{table}[!htb]
\centering
\caption[EarnHFT main results.]{Performance comparison on four cryptocurrency markets against six baselines (two policy-based, two value-based, two rule-based RL/rule methods). Arrows indicate the preferred direction; \fst{pink} marks the best entry per market and metric. TR, AVOL, and MDD are in percent.}
\label{tab:earnhft_main}
\renewcommand{\arraystretch}{1.15}
\resizebox{\textwidth}{!}{
\begin{tabular}{llcccccc}
\toprule
Market & Model & TR$\uparrow$ & ASR$\uparrow$ & ACR$\uparrow$ & ASoR$\uparrow$ & AVOL$\downarrow$ & MDD$\downarrow$ \\
\midrule
\multirow{7}{*}{BTCU}
 & DRA    & $-4.56$ & $-4.28$ & $-19.57$ & $-4.65$ & $42.25$ & $9.24$ \\
 & PPO    & $-3.61$ & $-5.25$ & $-22.74$ & $-5.71$ & $27.76$ & $6.41$ \\
 & CDQNRP & $-2.83$ & $-2.91$ & $-14.85$ & $-3.31$ & $37.61$ & $7.38$ \\
 & DQN    & $-3.48$ & $-12.37$ & $-35.01$ & $-11.86$ & \fst{$11.57$} & $4.09$ \\
 & MACD   & $-6.07$ & $-10.11$ & $-25.17$ & $-8.05$ & $24.84$ & $9.98$ \\
 & IV     & $-2.99$ & $-3.78$ & $-14.24$ & $-3.27$ & $31.35$ & $8.32$ \\
 & EarnHFT & \fst{$0.72$} & \fst{$1.22$} & \fst{$10.77$} & \fst{$0.93$} & $27.08$ & \fst{$3.07$} \\
\midrule
\multirow{7}{*}{BTCT}
 & DRA    & $-2.65$ & $-4.82$ & $-17.48$ & $-4.77$ & $21.18$ & $5.84$ \\
 & PPO    & $-0.60$ & $-14.74$ & $-35.80$ & $-0.14$ & $1.59$ & $0.65$ \\
 & CDQNRP & $-0.60$ & $-19.52$ & $-37.88$ & $-0.74$ & \fst{$1.20$} & \fst{$0.61$} \\
 & DQN    & $0.47$ & \fst{$4.21$} & \fst{$27.94$} & \fst{$1.14$} & $4.38$ & $0.66$ \\
 & MACD   & $-4.02$ & $-5.80$ & $-24.21$ & $-4.32$ & $26.87$ & $6.44$ \\
 & IV     & $-12.01$ & $-17.83$ & $-38.99$ & $-13.90$ & $27.68$ & $12.66$ \\
 & EarnHFT & \fst{$0.99$} & $1.34$ & $7.76$ & $1.01$ & $32.40$ & $5.61$ \\
\midrule
\multirow{7}{*}{ETH}
 & DRA    & $-33.37$ & $-9.06$ & $-32.23$ & $-9.20$ & $163.25$ & $45.88$ \\
 & PPO    & $-22.61$ & $-10.11$ & $-31.17$ & $-10.39$ & $96.12$ & $31.17$ \\
 & CDQNRP & $-6.82$ & $-24.41$ & $-40.19$ & $-3.11$ & \fst{$11.46$} & \fst{$6.96$} \\
 & DQN    & $-11.02$ & $-9.47$ & $-32.81$ & $-8.43$ & $47.76$ & $13.79$ \\
 & MACD   & $-4.29$ & $-1.78$ & $-8.71$ & $-1.19$ & $79.64$ & $16.35$ \\
 & IV     & $-27.42$ & $-12.27$ & $-36.01$ & $-9.00$ & $99.67$ & $33.96$ \\
 & EarnHFT & \fst{$4.52$} & \fst{$2.92$} & \fst{$14.30$} & \fst{$1.78$} & $67.92$ & $13.89$ \\
\midrule
\multirow{7}{*}{GALA}
 & DRA    & $10.56$ & $4.77$ & $41.63$ & $4.44$ & $92.43$ & $10.60$ \\
 & PPO    & $10.56$ & $4.77$ & $41.63$ & $4.44$ & $92.43$ & $10.60$ \\
 & CDQNRP & $5.22$ & $4.51$ & $39.42$ & $4.16$ & $47.27$ & $5.41$ \\
 & DQN    & $2.94$ & $3.55$ & $32.02$ & $2.66$ & \fst{$34.08$} & \fst{$3.78$} \\
 & MACD   & $2.37$ & $1.79$ & $11.45$ & $1.22$ & $62.89$ & $9.84$ \\
 & IV     & $13.95$ & $6.74$ & $55.67$ & $5.41$ & $81.79$ & $9.91$ \\
 & EarnHFT & \fst{$19.41$} & \fst{$9.77$} & \fst{$79.08$} & \fst{$7.79$} & $74.94$ & $9.26$ \\
\bottomrule
\end{tabular}}
\end{table}

\begin{figure}[!htb]
\centering
\includegraphics[width=0.6\textwidth]{Chapter4/earnhft_trading_example}
\caption[An EarnHFT trading episode.]{A representative EarnHFT trading episode on ETH. The agent opens and closes a position within roughly thirty seconds, capturing a movement that registers only as a pullback at the minute scale.}
\label{fig:earnhft_trading}
\end{figure}

\subsection{Effectiveness of the Hierarchical Framework}
\label{sec:earnhft_hierarchy}

\fref{fig:earnhft_hierarchy} compares the router with each specialist in its pool. Bull and rally specialists tend to buy and hold and excel in the bull market (GALA); sideways specialists trade little and hold; pullback and bear specialists tend to close positions and excel in the bear market (ETH). The router combines these strengths and delivers the best profit on all four datasets. \fref{fig:earnhft_selection} shows the router's selection distribution: high-volatility datasets (ETH, GALA) elicit a balanced distribution across trends because dynamics change frequently, whereas low-volatility datasets (BTCT, BTCU) concentrate on a few dynamics.

\begin{figure}[!htb]
\centering
\includegraphics[width=0.85\textwidth]{Chapter4/earnhft_hierarchical_comparison}
\caption[Router versus specialist pool.]{Comparison of the router against each specialist in its pool across the four datasets. Each specialist excels in the trend it was trained for; the router combines their strengths.}
\label{fig:earnhft_hierarchy}
\end{figure}

\begin{figure}[!htb]
\centering
\includegraphics[width=0.7\textwidth]{Chapter4/earnhft_trend_selection}
\caption[Router selection distribution.]{Distribution of the router's specialist selections across the four datasets. High-volatility markets elicit a more balanced distribution; low-volatility markets concentrate on fewer dynamics.}
\label{fig:earnhft_selection}
\end{figure}

\subsection{Effectiveness of the Optimal Action Value}
\label{sec:earnhft_ablation}

\tref{tab:earnhft_ablation} ablates the optimal value supervisor (OS) and optimal actor (OA) on ETH and GALA, measuring the steps to converge (CS), the converged reward sum (RS), and the average holding length (AHL). On GALA, adding OS to plain DDQN reduces the steps to converge to about 15\% while raising the return; OA further raises the return at the cost of more steps. On ETH, a bull market, OS does not reduce CS but substantially raises the return. OS is the more effective component because it supplies information about all actions and adapts as the policy changes, whereas OA only provides demonstrations.

\begin{table}[!htb]
\setlength\tabcolsep{5pt}
\centering
\caption[Ablation of the Q-teacher.]{Ablation of the optimal value supervisor (OS) and optimal actor (OA) on ETH and GALA. CS is steps to converge, RS the converged reward sum, AHL the average holding length.}
\label{tab:earnhft_ablation}
\begin{tabular}{cccccccc}
\toprule
\multirow{2}{*}{OA} & \multirow{2}{*}{OS} & \multicolumn{3}{c}{GALA} & \multicolumn{3}{c}{ETH} \\
\cmidrule(lr){3-5}\cmidrule(lr){6-8}
 & & CS$\downarrow$ & RS$\uparrow$ & AHL & CS$\downarrow$ & RS$\uparrow$ & AHL \\
\midrule
\checkmark & \checkmark & $78848$ & \fst{$4.43$} & $448$ & $102400$ & \fst{$12.32$} & $81.3$ \\
\checkmark &            & $102400$ & $0.24$ & $38.7$ & $102400$ & $-1.40$ & $4.15$ \\
           & \checkmark & \fst{$4608$} & $2.89$ & $147$ & \fst{$30720$} & $4.87$ & $35.8$ \\
           &            & $30720$ & $-0.01$ & $284$ & $30720$ & $-29.6$ & $39.1$ \\
\bottomrule
\end{tabular}
\end{table}

\subsection{Trading Behaviour and Computational Cost}
\label{sec:earnhft_behaviour}

The trading behaviour of EarnHFT varies with the underlying market in a manner consistent with the unit cost of each asset and the stability of its trend. On the two Bitcoin datasets (BTCT and BTCU), the high unit price makes each round trip relatively expensive in dollar terms, and the agent trades less frequently and holds positions only when the price excursion is comfortably larger than the round-trip commission. On ETH, where the trend reverses sharply within the test window, the agent trades much more often to capture the short pullbacks that the minute-level router exposes. On GALA, which trends upward through most of the test period without major reversals, the agent holds positions for long stretches and lets the trend do the work. In aggregate, EarnHFT trades less frequently in stable bull markets and more frequently in volatile or bear markets---behaviour that matches what a discretionary high-frequency trader would do, but emerges from the hierarchy without any hand-coded regime detector.

\paragraph{Computational cost.} All experiments run on a single NVIDIA RTX~4090 GPU. The full pipeline---building the Q-teacher table, training the $200$ specialist agents under the four preference values $\beta\in\{-90,-10,30,100\}$, segmenting and labelling the validation set, selecting the pool, and training the minute-level DDQN router---completes in approximately ten hours for each market, dominated by stage II specialist training. At inference time the only on-policy work is one specialist DDQN forward pass per second and one router forward pass per minute, both of which fit comfortably inside the second-level decision budget on a commodity GPU.

%----------------------------------------------------------------------------------------
\section{Chapter Summary}
\label{sec:earnhft_summary}

This chapter presented EarnHFT, a three-stage hierarchical RL framework that makes high-frequency trading tractable for RL. It addresses the extreme horizon with a dynamic-programming Q-teacher that, under the exogenous-price assumption of \sref{sec:earnhft_stage1} and a fixed historical training window, preserves the optimal Bellman fixed point of the supervised update while accelerating learning empirically; and it addresses regime non-stationarity by training a pool of trend specialists and a minute-level router that switches among them. EarnHFT outperforms six strong baselines across four markets, exceeding the runner-up by at least 30\% in profitability, and the ablations confirm the contribution of each component. In the language of the thesis, EarnHFT is the first instantiation of the decompose--specialise--select pattern, and its selector---a DDQN router trained by reinforcement learning against the trading reward---is the most tightly coupled of the three. The next chapter retains the pattern but confronts a setting, leveraged futures, where the router must additionally manage risk and the agent must become aware of its own competence boundary.
